\section{What the algorithm gives beyond a yes-or-no answer}
\label{sec:algorithm-consequences}

The consistency construction yields more than a gate-only exponent.
It identifies which part of a supplied circuit makes counting expensive,
and its exact counts can guide the construction of individual solutions.
These are consequences of the proved algorithm, using established
conditioning and self-reduction ideas. They are not claims of historical
priority for search or sampling from counts.

\subsection{The excess above a tree is a useful parameter}

Let $s$ be the supplied gate budget, $u$ the declared input count, and
$b$ the inputs left in a nontrivial normalized output cone. Define
\[
 d=s+1-b.
\]
The input-count lemma ensures $d\ge0$. The graph proof gives
$\kappa\le d$: every used input consumes one unit of the same budget
that would otherwise permit a consistency cycle. This is a statement
about the supplied representation, not its minimum possible size.

\begin{corollary}[Counting with logarithmic excess]
\label{cor:low-excess}
For each fixed $\delta>0$, exact counting takes time
\[
 \poly_\delta(s+u+1)\,2^{(1/3+\delta)d}.
\]
In particular, for every fixed $c\ge0$, circuits with
$d\le c\log(s+u+1)$ can be counted in polynomial time. The same
conclusion holds whenever $\kappa=O(\log(s+u+1))$, even if the
gate-minus-input estimate $d$ is larger.
\end{corollary}
\begin{proof}
The proof of Corollary~\ref{cor:gate-counting} first obtains time
$\poly_\delta(s+u+1)2^{(1/3+\delta)\kappa}$ and then bounds
$\kappa$ by $d$. If $d\le c\log(s+u+1)$, the exponential factor
is at most $(s+u+1)^{c(1/3+\delta)}$. Fixing, for example,
$\delta=1/6$ makes the total a polynomial of fixed degree.
The same calculation applies directly to $\kappa$.
\end{proof}

A binary AND tree on $u$ inputs has $s=u-1$, $b=u$, and $d=0$.
Its single accepting assignment is easy to count even though there
are $2^u$ candidate assignments. The explanation is not that few
solutions make counting easy: it is that the circuit has no excess
consistency to maintain. The running three-gate example has $b=3$
and $d=1$; its entire network also reduces by the elementary table
operations. Neither example needs the asymptotic graph theorem.

\paragraph{A sharper reading of the four-gates threshold.}
The $1/4$ coefficient protects against every possible value of $b$.
If all declared inputs survive normalization, $b=u$, the more precise
bound has limiting input coefficient $(c-1)/3$ at gate budget
$s\le cu$, provided $1\le c<4$. In contrast, the gate-only
coefficient is $c/4$. For example:
\begin{center}
\begin{tabular}{@{}lcc@{}}
\toprule
Gate budget & Arbitrary used-input count & All $u$ inputs remain used\\
\midrule
$s\le2u$ & $1/2$ & $1/3$\\
$s\le3u$ & $3/4$ & $2/3$\\
$s\le(7/2)u$ & $7/8$ & $5/6$\\
\bottomrule
\end{tabular}
\end{center}
The entries are limiting coefficients in an exponent of $2^u$.
Every entry requires arbitrarily small fixed positive slack and a
polynomial prefactor depending on that slack; these are not timing
measurements or endpoint bounds. Substituting $b=u$ into the proved
formula gives the right column directly.

More generally, with $r=b/u$ and $c=s/u$, the limiting envelope is
\[
 \min\{r,(c-r)/3\},
\]
up to fixed positive slack and the additive gate-budget constant.
Increasing $b$ makes enumeration more expensive but leaves fewer
cycles for contraction. The two costs balance near $b=(s+1)/4$.
Thus the most expensive value in this guarantee need not be the
largest number of used inputs. This is a balance of upper bounds,
not an example attaining either running-time bound or a hardness
theorem. Unused inputs merely multiply the exact answer by a power
of two; writing that answer costs polynomially many bits.

\subsection{Conditioning without losing the original budget}

Suppose some input values have already been chosen. We want to count
how many accepting assignments extend those choices. Re-normalizing
the restricted circuit changes $b$, so substituting its new parameters
into $s+1-b$ does not itself prove preservation of the original bound.
There is a simpler way: keep the original graph and change its tables.

Preserving only the unconditioned total would not justify this step.
The predicates $x\land\neg z$ and $\neg x\land z$ each have one
model on the named universe $(x,z)$, but pinning $x=1$ leaves one
model of the first and none of the second. We need the stronger
invariant that every boundary-conditioned count is preserved.

\begin{lemma}[Pinning preserves the counting layout]
\label{lem:pinning-layout}
Fix a nontrivial normalized output cone and its original consistency
graph of rank $\kappa$. For any partial assignment $e$ to the $u$
declared inputs, the number $Z(e)$ of satisfying assignments extending
$e$ can be computed in time
\begin{equation}\label{eq:pinned-counting}
 T_\delta=\poly_\delta(s+u+1)\,
 2^{\min\{b,(1/3+\delta)\kappa\}},
 \qquad \delta>0\text{ fixed},
\end{equation}
using the original graph reduction and vertex layout for every query.
Space is bounded by a polynomial times the same exponential factor.
Constant and designated-input outputs can be handled directly.
\end{lemma}
\begin{proof}
At a used input fixed by $e$, replace its equality table by the table
that is one precisely when all incident values equal the prescribed
bit. No vertex, edge, or index is added. When a high-degree equality
table is expanded as a ternary tree, impose the prescribed bit at one
of its ternary tables. The other equality tables propagate that bit
through the tree. A consistent external assignment still has exactly
one internal extension, and an inconsistent one has none.

The low-degree, loop, and parallel-edge reductions depend only on the
graph and are exact integer identities for the modified tables as
well. The same residual graph and layout therefore give the same
frontier sizes and operation bound. All entries remain counts of
assignments to $O(s+1)$ internal indices, with polynomial bit length.
For each originally unused input not fixed by $e$, multiply the
answer by two; pinned unused inputs contribute one instead. This
factor has at most $u+1$ bits.

Alternatively enumerate the original $b$ used input bits, discarding
assignments inconsistent with $e$. This costs at most a polynomial
times $2^b$. Choose between the two original guarantees. The graph
and layout can be prepared once, while the integer tables are
recomputed for each query. Storing queries sequentially suffices for
the stated space bound.
\end{proof}

\subsection{Search, rank, and exact uniform generation}

Order the declared input bits once and use lexicographic order on
their assignments, with $0<1$. A prefix $p$ has two disjoint sets
of extensions, so its exact counts satisfy
\[
 Z(p)=Z(p0)+Z(p1).
\]
This identity is the bridge from counting all solutions to selecting
one. Count-guided construction and generation are classical; see
\citet[Theorem 3.3 and the discussion of fair coins in Section 2]
{JerrumValiantVazirani1986}. The point here is that pinning retains
the original circuit's numerical bound.

\begin{corollary}[Constructing and indexing satisfying assignments]
\label{cor:search-unrank}
Within a polynomial factor of \eqref{eq:pinned-counting}, one can
deterministically find a satisfying assignment or report that none
exists; compute the zero-based lexicographic rank of a supplied
satisfying assignment; or reconstruct the satisfying assignment of
any supplied rank $R$ with $0\le R<Z(\varnothing)$.
The same gate-only coefficient approaching $1/4$ consequently
applies to all three tasks.
\end{corollary}
\begin{proof}
First compute $Z=Z(\varnothing)$. If $Z=0$, report unsatisfiability.
Otherwise, maintain a prefix $p$, an integer $R$, and the invariant
$0\le R<Z(p)$. Compute $A=Z(p0)$. If $R<A$, append zero; otherwise
append one and replace $R$ by $R-A$. The identity above preserves
the invariant. After $u$ steps the prefix is the unique satisfying
assignment of the original rank. Taking $R=0$ gives the first
satisfying assignment.

Conversely, verify a supplied assignment by direct circuit evaluation.
Along its prefixes, add $Z(p0)$ to an accumulator whenever its next
bit is one. These are exactly the disjoint lexicographically earlier
branches, so their sum is the rank. Both procedures use at most $u+1$
count queries and arithmetic on $O(u+1)$-bit integers. Sequential
queries multiply time by a polynomial and do not multiply space.
Finally use $\kappa\le s+1-b$ and the same balancing calculation
as in Corollary~\ref{cor:gate-counting}.
\end{proof}

For the running example the accepting assignments are
$010,101,110$, in that order. The zero branch for the first bit has
one accepting extension and the one branch has two. Unranking $R=2$
therefore selects the one branch and changes the residual rank to
one; continuing in the same way selects $110$. Choosing each
nonempty branch with probability one half would give the first
assignment probability $1/2$, rather than $1/3$.

\begin{corollary}[Exact uniform sampling with expected time]
\label{cor:exact-sampling}
For a satisfiable circuit, an exactly uniform satisfying assignment
can be sampled with expected running time bounded by a polynomial
times the exponential factor in \eqref{eq:pinned-counting}, using
independent fair random bits. Unsatisfiability is reported
deterministically within the same counting bound.
\end{corollary}
\begin{proof}
Compute $Z$. If $Z=0$, report unsatisfiability. If $Z=1$, unrank zero. For $Z>1$, put
$L=\lceil\log Z\rceil$. Draw $L$ independent fair bits, interpreted
as an integer $R\in\{0,\ldots,2^L-1\}$. Reject and repeat when
$R\ge Z$. Each trial accepts with probability $Z/2^L>1/2$, so
fewer than two trials are needed in expectation. Conditional on
acceptance, every rank is equally likely. Unranking is a bijection,
so every satisfying assignment has probability exactly $1/Z$.
Since $L\le u$, the expected random-bit and integer-processing
cost is polynomial, and Corollary~\ref{cor:search-unrank} bounds the
deterministic remainder.
\end{proof}

The sampler returns a solution with probability one and in finite expected
time; a measure-zero random stream can reject forever. The expected-time
qualifier is essential for this fair-bit implementation. With a fixed worst-case bound of $L$ coin
flips, every output probability is an integer multiple of $2^{-L}$,
even if some branches stop early. It cannot be exactly $1/3$ on
each of three solutions. Allowing rejection, or allowing a declared
failure outcome, resolves this elementary obstruction; silently
assuming a unit-cost arbitrary biased coin does not.

Finally, a uniform satisfying \emph{full assignment} need not give
a uniform positive \emph{projected input}. In the running example,
projecting the uniform distribution on $010,101,110$ onto $x$ gives
probabilities $1/3$ and $2/3$. Uniformity over witnesses weights each
ordinary input by its number of witnesses. The algorithm here counts
and samples full satisfying assignments; counting or sampling distinct
projected inputs is a different task.
